% !TEX root = ../main.tex
% ======================================================================
% Section 3 – Rappel de la problématique et des objectifs
% Grille : « Cadrage & problématique » – 6 pts (partagés avec la section 2).
% Longueur indicative : ~1 à 1,5 page.
% ======================================================================
% TODO : ce que \fedesante{} visait à résoudre (ex. « entraîner un modèle performant sans centraliser les dossiers patients »).
% TODO : formuler des objectifs SMART (spécifiques, mesurables, atteignables, réalistes, temporellement définis).
% TODO : énoncer le mandat de l'équipe : analyser de manière critique si la promesse suffit, et quels enjeux éthiques, légaux et de gouvernance subsistent.

\subsection{Problématique visée par \fedesante{}}
\label{sec:problematique-visee}
\arediger{problématique}

\subsection{Objectifs SMART}
\label{sec:objectifs-smart}
\arediger{objectifs}
% Suggestion : un tableau booktabs Objectif | Indicateur | Cible | Échéance.

\subsection{Mandat de l'analyse critique}
\label{sec:mandat}
\arediger{mandat}
